\chapter{Business Valuation with Large Language Models}
\label{ch:business-valuation}

\begin{remark}[Learning Objectives]
After completing this chapter, the reader should be able to:
\begin{enumerate}
  \item Explain the core mechanics of Discounted Cash Flow (DCF) valuation and
        identify the inputs where LLMs can augment human analysis.
  \item Describe how machine learning models improve earnings and cash-flow
        forecasts relative to analyst consensus.
  \item Apply NLP-based and LLM-based peer-selection methods for comparable-company
        valuation and articulate their advantages over industry-code approaches.
  \item Use an LLM to generate and structure scenario narratives for stress-testing
        a DCF model.
  \item Critically evaluate the limitations of LLM-generated valuations, including
        hallucination risk and the difficulty of auditing chain-of-thought reasoning
        in financial contexts.
\end{enumerate}
\end{remark}

% ============================================================
\section{Introduction: Why Valuation Is Hard and Where LLMs Help}
\label{sec:valuation-intro}
% ============================================================

Business valuation sits at the intersection of quantitative modelling and
qualitative judgement.  A practitioner performing a DCF analysis must (i) forecast
future cash flows, (ii) select an appropriate discount rate, (iii) choose peer
firms for multiple-based cross-validation, and (iv) articulate scenario narratives
for bull and bear cases.  Steps (i) and (iii) involve pattern recognition across
large volumes of financial text --- annual reports, earnings-call transcripts,
analyst research notes, and macroeconomic commentary --- that is precisely the kind
of task at which large language models excel.

The past few years have seen a rapid expansion of academic work quantifying the
extent to which LLMs and their statistical predecessors add value in each of these
subtasks.  \citet{kim2024financial} demonstrate that GPT-4, presented with
standardised, anonymised financial statements, outperforms the median sell-side
analyst in predicting the \emph{direction} of next-year earnings changes --- a
finding that has immediate implications for the DCF forecasting step.  At the same
time, \citet{lopezlira2023chatgpt} show that ChatGPT's sentiment scores for news
headlines contain information about subsequent daily stock returns that is
orthogonal to classical dictionary-based sentiment measures, suggesting that LLMs
surface relevant value-relevant signals beyond what traditional financial NLP
captures.

This chapter is structured as follows.  Section~\ref{sec:dcf-foundations}
reviews DCF mechanics and situates LLM assistance within the modelling
workflow.  Section~\ref{sec:cf-forecasting} covers machine-learning and
LLM-based cash-flow forecasting.  Section~\ref{sec:peer-selection} examines
NLP-based comparable-company selection.  Section~\ref{sec:scenario-analysis}
introduces LLM-generated scenario analysis.  Section~\ref{sec:equity-research}
surveys the automation of equity research workflows.
Section~\ref{sec:valuation-limitations} discusses limitations and best practices.

% ============================================================
\section{DCF Valuation: Mechanics and LLM Touch-Points}
\label{sec:dcf-foundations}
% ============================================================

\subsection{The Core DCF Identity}
\label{sec:dcf-identity}

The enterprise value of a firm under the standard DCF framework is the present
value of its expected free cash flows (FCF) plus a continuing (terminal) value:

\begin{equation}
  V_0^{\mathrm{EV}}
  = \sum_{t=1}^{T} \frac{\mathrm{FCF}_t}{(1 + r_{\mathrm{WACC}})^t}
  + \frac{\mathrm{TV}}{(1 + r_{\mathrm{WACC}})^T},
  \label{eq:dcf-ev}
\end{equation}

where $r_{\mathrm{WACC}}$ is the weighted average cost of capital and
$\mathrm{TV}$ is the terminal value, typically computed via the Gordon Growth
Model:

\begin{equation}
  \mathrm{TV} = \frac{\mathrm{FCF}_{T+1}}{r_{\mathrm{WACC}} - g},
  \label{eq:gordon-tv}
\end{equation}

with $g$ the long-run nominal growth rate.  Free cash flow is obtained from
accounting statements as

\begin{equation}
  \mathrm{FCF}_t = \mathrm{EBIT}_t (1 - \tau)
    + \mathrm{D\&A}_t
    - \Delta\mathrm{NWC}_t
    - \mathrm{CapEx}_t,
  \label{eq:fcf}
\end{equation}

where $\tau$ is the marginal corporate tax rate, $\mathrm{D\&A}$ denotes
depreciation and amortisation, $\Delta\mathrm{NWC}$ is the change in net working
capital, and $\mathrm{CapEx}$ is capital expenditure.

\begin{remark}
  The Gordon Growth Model in \eqref{eq:gordon-tv} is valid only when
  $r_{\mathrm{WACC}} > g$.  In low-interest-rate environments this constraint
  can bind for high-growth technology firms, requiring alternative terminal-value
  approaches (exit multiples, H-model) that are discussed in
  Section~\ref{sec:scenario-analysis}.
\end{remark}

\subsection{LLM Touch-Points in the DCF Workflow}
\label{sec:dcf-touchpoints}

Table~\ref{tab:dcf-touchpoints} maps the four main DCF inputs to the specific
NLP or LLM capabilities that can augment each one.

\begin{table}[htbp]
\centering
\caption{LLM assistance in the DCF workflow.  Each row identifies a DCF input,
the information source, and the LLM capability that assists.}
\label{tab:dcf-touchpoints}
\small
\begin{tabular}{p{3.2cm} p{3.2cm} p{5cm}}
\toprule
\textbf{DCF Input} & \textbf{Primary Source} & \textbf{LLM Capability} \\
\midrule
FCF forecasts
  & 10-K/Q filings, earnings calls
  & Financial statement analysis~\cite{kim2024financial};
    earnings change prediction~\cite{chen2022predicting} \\[4pt]
Discount rate ($r_{\mathrm{WACC}}$)
  & Capital structure, market data
  & Risk-factor extraction from MD\&A; credit-quality scoring \\[4pt]
Peer firm selection
  & Industry codes, analyst reports
  & NER-based peer identification~\cite{covas2023ner};
    ML relative valuation~\cite{geertsema2023relative} \\[4pt]
Scenario narratives
  & Management guidance, macro data
  & LLM scenario generation~\cite{aldridge2024llm};
    Monte Carlo augmentation \\
\bottomrule
\end{tabular}
\end{table}

The table makes clear that LLMs are most naturally applied to the \emph{qualitative}
information synthesis steps (narrative generation, peer identification from business
descriptions) while statistical machine learning models address the \emph{quantitative}
forecasting steps (earnings changes, cash-flow paths).  In practice, the most
effective pipelines combine both: an LLM extracts structured signals from unstructured
text, and a gradient-boosted model or neural network converts those signals into
numerical forecasts.

% ============================================================
\section{Machine-Learning Cash-Flow and Earnings Forecasting}
\label{sec:cf-forecasting}
% ============================================================

\subsection{Why Analyst Forecasts Are Systematically Biased}
\label{sec:analyst-bias}

A long-standing empirical regularity in accounting research is that sell-side
analysts exhibit systematic optimism: consensus EPS forecasts are, on average,
too high, particularly at longer horizons.  The source of this bias is
multi-faceted --- incentive conflicts with investment banking, loss aversion in
career concerns, and the well-documented tendency to anchor on prior-period
guidance.  Machine-learning models trained on historical financial statement
data are immune to these behavioural distortions by construction, which
partially explains their out-of-sample advantage.

\begin{example}[Machine Learning Earnings Forecasting]
\label{ex:ml-earnings}
\citet{chen2022predicting} train gradient-boosted and random forest models on
a panel of U.S.~public firms spanning 1995--2019, using over 190 financial
statement line items as features.  Their best model reduces mean absolute
forecast error by roughly 15\% relative to the random-walk benchmark and by
12\% relative to consensus analyst forecasts in the one-year-ahead horizon.
Critically, the improvement concentrates in smaller firms, firms with high
earnings volatility, and firms with large accruals --- precisely the cases
where human analysts tend to struggle and where the DCF valuation is most
sensitive to forecast accuracy.
\end{example}

\subsection{LLMs for Financial Statement Analysis}
\label{sec:llm-fs-analysis}

The most influential recent result in the intersection of LLMs and valuation
is the finding by \citet{kim2024financial} that GPT-4 can perform financial
statement analysis at a level competitive with, and in some dimensions
exceeding, trained human analysts.  In their experiment, standardised and
anonymised balance sheets, income statements, and cash-flow statements are
provided to GPT-4 without any industry or firm identifiers.  The model is
asked to predict whether earnings will increase or decrease in the following
year.  GPT-4 achieves an accuracy of approximately 60\% on this task,
compared with 56\% for the median analyst in the same prediction window,
and 53\% for a simple random-walk baseline.  The LLM's relative advantage
is largest in periods with high analyst disagreement --- a signal that the
model extracts information that analysts collectively miss or discount.

\begin{remark}
  The \citet{kim2024financial} result does not imply that LLMs should replace
  analysts in full-cycle valuation.  The experiment strips away narrative
  context, industry knowledge, and management relationships that experienced
  analysts bring to a DCF.  Rather, the finding motivates using LLMs as a
  \emph{first-pass screening} tool to flag discrepancies between model
  predictions and analyst consensus before committing to a full valuation.
\end{remark}

\subsection{LLMs as Financial Analysts}
\label{sec:llm-as-analysts}

Beyond financial-statement analysis, a growing body of work examines whether
LLMs can replicate the full workflow of a sell-side analyst.
\citet{alonso2024llm} evaluate GPT-4o, Claude 3.5 Sonnet, and Gemini Advanced
on a structured financial analysis task covering revenue trends, profitability,
cash-flow management, and forward guidance for five major technology companies.
Their comparative study finds that the models differ systematically in their
analysis depth: GPT-4o produces the most granular quantitative breakdown, while
Claude 3.5 Sonnet provides broader strategic framing.  These differences mirror
the documented variation in human analyst styles across brokerage houses and
point to a promising research agenda: understanding which model
``personality'' is best suited to which valuation task.

\subsection{Machine Learning for Earnings Forecasting at Scale}
\label{sec:ml-earnings-scale}

Complementing single-model studies, \citet{chattopadhyay2025ml} conduct a
large-scale horse race of machine-learning methods for earnings-per-share
forecasting using data from the U.S.\ and seventeen international markets.
Their headline finding is that XGBoost --- a gradient-boosted tree ensemble ---
dominates linear benchmarks and neural network alternatives in out-of-sample
accuracy.  Crucially, the XGBoost forecasts exhibit the highest ``value
relevance,'' meaning their residuals are correlated with subsequent abnormal
returns in a direction consistent with fundamental mispricing.  This bridges
the gap between accounting-based earnings forecasting and asset pricing: accurate
machine-learning earnings forecasts contain information that the market has not
yet fully incorporated, providing a theoretically grounded rationale for using
such forecasts as DCF inputs rather than consensus analyst estimates.

% ============================================================
\section{Peer Firm Selection and Comparable Company Analysis}
\label{sec:peer-selection}
% ============================================================

\subsection{Limitations of Standard Industry Classification}
\label{sec:sic-limitations}

The traditional approach to comparable company selection relies on Standard
Industrial Classification (SIC) codes or their GICS/NAICS successors.
These hierarchical codes assign each firm a primary industry category based on
its main revenue stream.  The approach has well-documented deficiencies for
valuation purposes: firms within the same four-digit SIC code can differ
enormously in their product mix, capital intensity, and growth profile; and
diversified conglomerates are assigned a single code that may describe only a
fraction of their operations.  As a result, multiples-based cross-validation of
a DCF often introduces large pricing errors attributable to poor peer matching.

\subsection{Machine Learning Relative Valuation}
\label{sec:ml-relative-valuation}

\citet{geertsema2023relative} propose a machine-learning framework for relative
valuation that replaces hard industry boundaries with \emph{data-driven} peer
weights.  Their key insight is that the valuation multiple predicted by a
gradient-boosted model can be algebraically decomposed into a weighted average
of comparable firm multiples, where the weights are functions of the model's
feature-importance scores.  In out-of-sample tests across 1985--2019, their
ML-predicted multiples for enterprise value, price-to-earnings, and price-to-book
substantially outperform industry-median multiples, with mean absolute percentage
errors roughly 20--30\% lower.  Moreover, long-short portfolios that buy
undervalued firms (actual multiple below ML prediction) and short overvalued
firms generate statistically significant abnormal returns, providing a joint
test of valuation accuracy and market efficiency.

\begin{example}[Decomposing ML Multiples into Peer Weights]
\label{ex:ml-peer-weights}
Let $q_i$ be the observed EV/EBITDA multiple for firm $i$ and let
$\hat{q}_i = f(\mathbf{x}_i)$ be the model prediction, where $\mathbf{x}_i$
contains profitability ratios, growth rates, capital efficiency, and leverage.
\citet{geertsema2023relative} show that for tree-based models, the prediction
can be written as
\[
  \hat{q}_i = \sum_{j \neq i} w_{ij} q_j + \text{intercept},
\]
where $w_{ij} \geq 0$ and $\sum_j w_{ij} = 1$.  Firms with large $w_{ij}$
are the model-implied comparables for firm $i$, regardless of whether they
share an industry code.
\end{example}

\subsection{NLP-Based Peer Identification}
\label{sec:nlp-peers}

An orthogonal but complementary approach uses the \emph{text} of business
descriptions, corporate filings, and website content to identify peers.
\citet{covas2023ner} demonstrate that GPT-3.5-turbo, prompted to extract
products and services from company descriptions scraped from Wikipedia, achieves
an F-score of approximately 85\% in correctly identifying comparable companies
for equity valuation.  This substantially outperforms a standard SpaCy NER
baseline, which achieves only 61\%, confirming that the richer contextual
understanding of LLMs translates into meaningfully better peer identification.
The practical implication for valuation practice is that LLM-based peer
identification can be automated as a preprocessing step, with a human analyst
reviewing the top-$k$ suggestions rather than constructing the peer group from
scratch.

\begin{remark}
  NLP-based peer selection is particularly valuable for private-company valuation,
  where SIC codes may be unreliable or unavailable, and where the analyst must
  rely on fragmentary information (pitch decks, LinkedIn profiles, press releases)
  to construct a peer group.  LLMs excel at synthesising heterogeneous unstructured
  sources into a structured comparables list.
\end{remark}

% ============================================================
\section{Scenario Analysis and Stress-Testing with LLMs}
\label{sec:scenario-analysis}
% ============================================================

\subsection{The Role of Scenarios in DCF Valuation}
\label{sec:scenario-role}

A point estimate of $V_0^{\mathrm{EV}}$ from \eqref{eq:dcf-ev} is almost always
misleading.  DCF valuation is acutely sensitive to the assumed long-run growth
rate $g$ and the discount rate $r_{\mathrm{WACC}}$: a 50-basis-point change in
either can shift enterprise value by 15--25\% for a typical mid-cap firm.
Practitioners therefore complement the base case with scenario analysis ---
assigning probabilities to alternative states of the world (bull, base, bear)
and computing the probability-weighted valuation.  The difficulty lies in
constructing internally consistent scenario narratives: a ``bear'' scenario
for a retail firm must simultaneously articulate plausible assumptions about
consumer spending, interest rates, supply-chain costs, and competitive entry.

\subsection{LLM-Generated Scenario Narratives}
\label{sec:llm-scenarios}

\citet{aldridge2024llm} propose embedding LLM-generated scenario descriptions
within a Temporal Fusion Transformer (TFT) framework as an alternative to
classical Monte Carlo simulation.  In their approach, an LLM is prompted to
generate $K$ qualitatively distinct macroeconomic scenarios (e.g., ``sustained
rate hike cycle'', ``stagflation with supply shock'', ``soft landing with AI
productivity boom'').  Each scenario is converted into a structured covariate
vector that conditions the TFT's attention mechanism, producing a scenario-specific
forward path for free cash flows.  The resulting distribution over valuation
outcomes is richer than a standard Monte Carlo simulation because the scenarios
are grounded in coherent narrative logic rather than drawn independently from
marginal distributions.

\begin{definition}[Scenario-Conditioned Valuation Distribution]
\label{def:scenario-val}
Let $\mathcal{S} = \{s_1, \ldots, s_K\}$ be a set of $K$ LLM-generated scenarios
with prior probabilities $\{p_k\}_{k=1}^K$, $\sum_k p_k = 1$.  For each scenario
$s_k$, let $V_k$ be the enterprise value computed from the scenario-specific FCF
path.  The scenario-weighted valuation is
\begin{equation}
  \bar{V} = \sum_{k=1}^{K} p_k V_k,
  \label{eq:scenario-val}
\end{equation}
and the scenario uncertainty is measured by
\begin{equation}
  \sigma_V^2 = \sum_{k=1}^{K} p_k (V_k - \bar{V})^2.
  \label{eq:scenario-var}
\end{equation}
The distribution $\{(p_k, V_k)\}$ provides a richer uncertainty quantification
than the point estimate $V_0^{\mathrm{EV}}$ in \eqref{eq:dcf-ev}.
\end{definition}

The approach of \citet{aldridge2024llm} is notable for its hybrid design: it
leverages the narrative coherence of LLMs for scenario generation while delegating
numerical time-series forecasting to a purpose-built deep learning architecture.
This division of labour addresses the well-known weakness of LLMs in precise
arithmetic while capitalising on their strength in qualitative reasoning.

\subsection{Sensitivity Analysis and the WACC}
\label{sec:sensitivity-wacc}

Even without LLM-generated narratives, scenario analysis can be formalised as
a sensitivity table over $(g, r_{\mathrm{WACC}})$.  Let $V(g, r)$ denote the
enterprise value as a function of the two key parameters.  A first-order
Taylor expansion around the base-case values $(g_0, r_0)$ gives

\begin{equation}
  V(g, r) \approx V_0
    + \frac{\partial V}{\partial g}\bigg|_{g_0, r_0} (g - g_0)
    + \frac{\partial V}{\partial r}\bigg|_{g_0, r_0} (r - r_0),
  \label{eq:sensitivity}
\end{equation}

where, from \eqref{eq:dcf-ev} and \eqref{eq:gordon-tv}, the partial derivatives
can be computed in closed form for the terminal-value component.  LLMs can assist
in choosing the range over which to perturb $g$ and $r$ by extracting management
guidance from earnings-call transcripts and identifying the scenarios that analysts
consider most salient.

% ============================================================
\section{Equity Research Automation}
\label{sec:equity-research}
% ============================================================

\subsection{The Anatomy of an Equity Research Report}
\label{sec:er-anatomy}

A standard sell-side equity research report contains: (i) a company overview,
(ii) a thesis statement and key investment considerations, (iii) a financial
model with forecasts, (iv) a valuation section (DCF, EV/EBITDA, P/E comparables),
and (v) risk factors.  Each of these components draws on different information
sources --- public filings, earnings calls, industry reports, macro data ---
and requires different analytical skills.  LLMs are best positioned to assist
with (i), (ii), and (v), where synthesising large amounts of text into concise,
readable prose is the primary task.

\subsection{LLM Stock Price Forecasting}
\label{sec:llm-stock-forecast}

\citet{lopezlira2023chatgpt} conduct a large-scale experiment in which ChatGPT
is asked to classify news headlines as \emph{good}, \emph{bad}, or \emph{irrelevant}
for a given stock, without any financial training data or explicit stock-price
information.  The resulting ``ChatGPT score'' is significantly positively correlated
with next-day abnormal returns: stocks in the top quintile of ChatGPT sentiment
outperform those in the bottom quintile by approximately 0.5\% per day in the
2021--2022 sample period.  Importantly, this predictive power is orthogonal to the
Loughran--McDonald dictionary sentiment score \citep{loughran2011liability},
suggesting that LLMs capture semantic content that is not reducible to pre-trained
financial word lists.

The Lopez-Lira and Tang finding has important implications for equity research
automation.  It suggests that an LLM-powered news-monitoring system could
provide a real-time sentiment signal that supplements the analyst's fundamental
view.  The practical challenge is that the signal's predictive horizon is daily,
making it relevant for short-horizon trading rather than the multi-year DCF
horizon that underpins most equity research.

\subsection{EDGAR and Textual Disclosure Analysis}
\label{sec:edgar-text}

A large share of the information relevant to equity valuation is embedded in
regulatory filings submitted to EDGAR (Electronic Data Gathering, Analysis, and
Retrieval).  The text of MD\&A sections, risk factors, and footnotes to financial
statements contains forward-looking statements, qualitative risk assessments, and
management commentary that is not captured in the reported numbers.

LLMs are particularly well suited to mining EDGAR text at scale.  For instance,
\citet{kim2024financial} use the structured financial tables from SEC 10-K
filings as their primary input, relying on GPT-4's ability to perform
chain-of-thought reasoning over multi-table numerical data.  A complementary
approach, implicit in the machine-learning literature \citep{chen2022predicting},
is to construct features from the textual portions of filings --- MD\&A
readability, forward-looking statement density, and tone measures --- and combine
them with accounting variables in a predictive model.

% ============================================================
\section{Limitations and Best Practices}
\label{sec:valuation-limitations}
% ============================================================

\subsection{Hallucination in Financial Contexts}
\label{sec:hallucination-val}

The most acute limitation of LLM-assisted valuation is hallucination: the model
confidently generating financial figures, citations, or ratios that have no basis
in the source documents.  In a DCF context, hallucinated revenue forecasts or
invented comparable-company multiples can propagate through the model and produce
a plausible-looking but fundamentally incorrect valuation.  Best practices to
mitigate this risk include:

\begin{enumerate}
  \item \textbf{Retrieval-augmented generation (RAG).}  Ground the LLM's
        responses in retrieved passages from verified source documents (e.g.,
        the most recent 10-K filing).  The model is constrained to synthesise
        rather than invent.
  \item \textbf{Structured output with source attribution.}  Require the LLM
        to output a JSON object containing each numerical claim alongside the
        page and line number from which it was extracted.  Claims without
        verifiable provenance are flagged for human review.
  \item \textbf{Cross-model consistency checking.}  Run the same prompt through
        two models and flag discrepancies in numerical outputs exceeding a
        threshold (e.g., $\pm 5\%$).
  \item \textbf{Statistical benchmarking.}  Compare LLM-derived earnings
        forecasts against a simple time-series baseline (e.g., AR(1) on EBITDA).
        LLM outputs that deviate materially from the statistical baseline should
        trigger additional scrutiny.
\end{enumerate}

\subsection{Auditability and Regulatory Considerations}
\label{sec:auditability}

Regulatory frameworks governing investment research (MiFID II in Europe, SEC
Regulation AC in the United States) impose certification requirements on
analyst forecasts.  A valuation report co-generated by an LLM raises the
question of \emph{who} is certifying the forecast: the human analyst who
issued the prompt, or the AI model that generated the content?  At the time of
writing, no jurisdiction has provided definitive guidance on this question.
Practitioners should maintain a complete audit trail of LLM interactions ---
prompts, model version, temperature settings, and retrieved context --- to
support any subsequent regulatory review.

\subsection{Data-Snooping and Evaluation Pitfalls}
\label{sec:data-snooping}

Machine-learning valuation models are susceptible to look-ahead bias and
overfitting to historical correlations that do not persist out of sample.
\citet{geertsema2023relative} address this by conducting a careful walk-forward
cross-validation study in which the model is retrained monthly using only data
available at each point in time.  Practitioners replicating such methods should
be equally rigorous: models trained on data through 2019 and evaluated on 2020
data (the COVID-19 shock) will exhibit large forecast errors irrespective of
model quality, because the distributional shift is unprecedented.

% ============================================================
\section{Summary}
\label{sec:valuation-summary}
% ============================================================

This chapter has surveyed the growing evidence that LLMs and machine-learning
methods can add measurable value at multiple stages of the business valuation
workflow.  Key takeaways are:

\begin{itemize}
  \item \textbf{Financial statement analysis:} GPT-4 outperforms the median
        analyst in predicting earnings direction from anonymised financial
        statements \citep{kim2024financial}, suggesting a role for LLMs in
        the FCF forecasting step.
  \item \textbf{Earnings forecasting:} Gradient-boosted models trained on
        detailed accounting variables reduce forecast error by 10--15\%
        relative to consensus \citep{chen2022predicting}, with XGBoost
        forecasts exhibiting significant value relevance internationally
        \citep{chattopadhyay2025ml}.
  \item \textbf{Peer selection:} ML-based relative valuation substantially
        outperforms SIC-code peer matching \citep{geertsema2023relative}, and
        GPT-based NER achieves 85\% F-score in identifying comparable
        companies from business descriptions \citep{covas2023ner}.
  \item \textbf{Scenario generation:} LLM-conditioned Temporal Fusion
        Transformers provide richer scenario distributions than independent
        Monte Carlo draws \citep{aldridge2024llm}.
  \item \textbf{Equity research:} LLM sentiment scores contain return-predictive
        information orthogonal to dictionary-based measures
        \citep{lopezlira2023chatgpt}, while multi-model comparisons reveal
        systematic differences in analysis depth \citep{alonso2024llm}.
\end{itemize}

These findings collectively support a ``human-in-the-loop'' view of LLM-assisted
valuation: models excel at information synthesis, first-pass screening, and
scenario structuring, while the final judgements about discount rates, terminal
growth assumptions, and the weighting of competing scenarios remain the
responsibility of the human analyst.  Chapter~6 extends these ideas to the
construction of full LLM-based agent pipelines that can execute a multi-step
valuation workflow autonomously, with human review at designated checkpoints.
